% ============================================================================
%  My Notes: Stacks  (CIS 4020)
%  Requires classnotes.sty in the same folder or your local texmf tree.
%  Compile twice with pdflatex so "Page X of Y" resolves.
% ============================================================================
\documentclass[11pt]{article}
\usepackage[margin=1in]{geometry}
\usepackage[pagenumbers]{classnotes}
\usepackage{tikz}
\usetikzlibrary{positioning,arrows.meta,calc}
\newcolumntype{L}[1]{>{\raggedright\arraybackslash}p{#1}}

\classnotesmeta{Stacks: Notes (CIS 4020)}

% ---- Java / C++ listing styles, built from the classnotes palette ----
\lstdefinestyle{java}{
  language=Java,
  backgroundcolor=\color{codebg},
  basicstyle=\ttfamily\small\color{navy},
  keywordstyle=\color{teal}\bfseries,
  commentstyle=\color{codecmt}\itshape,
  stringstyle=\color{amber},
  showstringspaces=false,
  breaklines=true, keepspaces=true, columns=fullflexible,
  frame=single, rulecolor=\color{rulecol},
  framesep=6pt, xleftmargin=6pt, xrightmargin=6pt, aboveskip=8pt, belowskip=8pt,
}
\lstdefinestyle{cpp}{
  language=C++,
  backgroundcolor=\color{codebg},
  basicstyle=\ttfamily\small\color{navy},
  keywordstyle=\color{teal}\bfseries,
  commentstyle=\color{codecmt}\itshape,
  stringstyle=\color{amber},
  showstringspaces=false,
  breaklines=true, keepspaces=true, columns=fullflexible,
  frame=single, rulecolor=\color{rulecol},
  framesep=6pt, xleftmargin=6pt, xrightmargin=6pt, aboveskip=8pt, belowskip=8pt,
}
\lstdefinestyle{pseudo}{
  backgroundcolor=\color{codebg},
  basicstyle=\ttfamily\small\color{navy},
  commentstyle=\color{codecmt}\itshape,
  showstringspaces=false,
  breaklines=true, keepspaces=true, columns=fullflexible,
  frame=single, rulecolor=\color{rulecol},
  framesep=6pt, xleftmargin=6pt, xrightmargin=6pt, aboveskip=8pt, belowskip=8pt,
}

\begin{document}

\classnotestitle{CIS 4020 \quad$\cdot$\quad DATA STRUCTURES}{Stacks}{Notes on the ADT, both implementations, and its classic applications, with Java and C/C++ examples}

\notebox{A note before you begin.}{A word before you read. These are my notes from years of teaching this unit, cleaned up so they are
readable. While C/C++ was my first language, I am a bit rusty because I have taught Java for years now.
For these notes, I have been using AI to help type these up, check for issues, and fill in explanations, which
means a wrong or imprecise statement can slip in and survive longer than it should, and I do not fully
trust the AI, so I am still editing. They are not meant to be a study guide, they are not gospel, and not
everything here shows up in a lecture (and some lecture items are not here). If you find something that
looks wrong, contradicts the book, or just does not behave the way I claim when you run it, tell me. You
will be doing students and me a favor, and you will understand it better for having caught it. The textbook
is the tiebreaker, and the compiler beats us both. We are a two-language course, thus I will do my best to
show everything in Java (CIS) and C/C++ (CS).
}

\section{The stack ADT}

A stack is an abstract data type, and the wall between what it offers and how it is built
behind the scenes works the same way it did for the linked list, covered in that unit's
notes. We will not repeat that discussion here, only apply it. A stack ADT specifies a
small set of operations, push, pop, peek (or top), isEmpty, and size, without saying
anything about how those operations are carried out underneath. This unit covers two ways
to carry them out: an array-based version and a linked version.

\section{Terminology}

\begin{itemize}[leftmargin=1.4em]
  \item \textbf{LIFO.} Last in, first out. The most recently added item is the first one
    removed.
  \item \textbf{Top.} The end of the stack where every push and pop happens. The other end
    of the stack, sometimes called the bottom, is never touched directly.
  \item \textbf{Push.} Adds a new item at the top.
  \item \textbf{Pop.} Removes and returns the item at the top.
  \item \textbf{Peek (or top).} Returns the item at the top without removing it.
  \item \textbf{Underflow.} An attempt to pop or peek an empty stack. Every implementation
    needs to guard against this.
  \item \textbf{Overflow.} An attempt to push onto a stack that has reached a fixed
    capacity. This only applies to an array-based stack that does not resize itself.
\end{itemize}

\begin{figure}[htbp]
\centering
\begin{tikzpicture}[
  cell/.style={draw=navy, thick, minimum height=0.65cm, minimum width=1.7cm, font=\ttfamily}
]
\node[cell, fill=white] (b1) {12};
\node[cell, fill=white, above=0 of b1] (b2) {47};
\node[cell, fill=white, above=0 of b2] (b3) {9};
\draw[-{Triangle[length=2.5mm,width=2mm]}, teal, thick] ($(b3.east)+(1.1,0)$) -- (b3.east)
  node[midway, above, font=\small\bfseries, teal] {top};
\node[below=0.15cm of b1, font=\small\itshape] {bottom of stack};
\end{tikzpicture}
\caption{A stack holding three items. Every push and pop happens at the top; the bottom of
the stack is never touched directly.}
\label{fig:stack-anatomy}
\end{figure}

\section{Array-based versus linked implementation}

An array-based stack keeps its items in a contiguous block and tracks the top with an
index. A linked stack keeps its items in nodes, the same node and link idea from the
linked list unit, and tracks the top with a reference to the first node in the chain.

\begin{center}
\begin{minipage}{\linewidth}
\centering
\captionof{table}{Array-based versus linked implementations of a stack.}
\renewcommand{\arraystretch}{1.3}
\begin{tabular}{>{\bfseries}l L{5.6cm} L{5.6cm}}
\toprule
\rowcolor{navy!8}
Aspect & Array-based & Linked \\
\midrule
Memory layout & Contiguous block, plus a top index & Scattered nodes, top is a reference to the first node \\
Capacity & Fixed, or resized (typically doubled) once full & Grows and shrinks one node at a time, no fixed limit \\
Push cost & O(1), occasionally O(n) when a resize is triggered & O(1) \\
Pop cost & O(1) & O(1) \\
Overflow & Possible if fixed-size and full & Not applicable, bounded only by available memory \\
\bottomrule
\end{tabular}
\end{minipage}
\end{center}

\notebox{A note on the code.}{The Java and C/C++ snippets below are not the exact code in
your textbook, and I expect the two to look different in places. I simplify and reorganize
for teaching purposes, and the book may name methods differently, split the class up
differently, or handle a case, such as resizing or an empty stack, in a different order.
These are snippets, not full listings, meant to show the key idea in each operation, not a
complete class you can compile as-is. When the versions genuinely disagree, the textbook
and the compiler settle it, not these notes.}

\section{Code snippets: array-based stack}

The array-based stack keeps an array of items and an integer, often called top, that marks
the index one past the last used slot (equivalently, the current size of the stack).

\subsection*{Java}

\begin{lstlisting}[style=java]
private Object[] items;
private int top;   // number of items currently on the stack

public void push(T value) {
    if (top == items.length) {
        items = Arrays.copyOf(items, items.length * 2);   // resize, doubling capacity
    }
    items[top] = value;
    top++;
}

@SuppressWarnings("unchecked")
public T pop() {
    if (isEmpty()) {
        throw new NoSuchElementException("stack is empty");
    }
    top--;
    T value = (T) items[top];
    items[top] = null;   // avoid holding a reference the stack no longer owns
    return value;
}
\end{lstlisting}

\subsection*{C/C++}

\begin{lstlisting}[style=cpp]
T* items;
int capacity;
int top;   // number of items currently on the stack

void push(const T& value) {
    if (top == capacity) {
        resize(capacity * 2);   // allocate a bigger array and copy items over
    }
    items[top] = value;
    top++;
}

T pop() {
    if (isEmpty()) {
        throw std::underflow_error("stack is empty");
    }
    top--;
    return items[top];
}
\end{lstlisting}

\section{Code snippets: linked stack}

The linked stack keeps a single reference, top, pointing at the first node in the chain.
Pushing and popping both happen at that one reference, so neither operation ever has to
walk the list.

\subsection*{Java}

\begin{lstlisting}[style=java]
private Node<T> top;

public void push(T value) {
    Node<T> newNode = new Node<>(value);
    newNode.setNext(top);   // new node points at the old top
    top = newNode;          // top now points at the new node
}

public T pop() {
    if (isEmpty()) {
        throw new NoSuchElementException("stack is empty");
    }
    T value = top.getData();
    top = top.getNext();
    return value;
}
\end{lstlisting}

\subsection*{C/C++}

\begin{lstlisting}[style=cpp]
Node<T>* top;

void push(const T& value) {
    Node<T>* newNode = new Node<T>(value);
    newNode->next = top;   // new node points at the old top
    top = newNode;         // top now points at the new node
}

T pop() {
    if (isEmpty()) {
        throw std::underflow_error("stack is empty");
    }
    Node<T>* oldTop = top;
    T value = oldTop->data;
    top = top->next;
    delete oldTop;          // free the node we just unlinked
    return value;
}
\end{lstlisting}

\section{Applications, in pseudocode}

Both classic applications below only need push, pop, peek, and isEmpty, which is part of
why the stack ADT shows up so often once you know to look for it.

\subsection*{Checking for balanced parentheses}

\begin{lstlisting}[style=pseudo]
create an empty stack
for each character c in the input, left to right:
    if c is an opening symbol ( ( [ { ):
        push c
    else if c is a closing symbol ( ) ] } ):
        if the stack is empty:
            return "not balanced"
        pop the top symbol and check that it matches c
        if it does not match:
            return "not balanced"
after the loop, if the stack is empty:
    return "balanced"
else:
    return "not balanced"   // some opening symbol was never closed
\end{lstlisting}

\subsection*{Evaluating a postfix expression}

Postfix notation (also called reverse Polish notation) writes an operator after its two
operands, for example \texttt{3 4 +} instead of \texttt{3 + 4}, which avoids any need for
parentheses or operator precedence rules.

\begin{lstlisting}[style=pseudo]
create an empty stack
for each token t in the postfix expression, left to right:
    if t is a number:
        push t
    else if t is an operator (+, -, *, /):
        pop the top value into b
        pop the top value into a
        compute result = a t b     // for example, a + b if t is "+"
        push result
after the loop, pop and return the single remaining value on the stack
\end{lstlisting}

\section{Stacks and recursion}

The call stack that a program uses to keep track of recursive calls, from the first unit
of this course, is a stack in exactly this sense: each call pushes a new frame holding that
call's local variables and its return address, and returning from a call pops the
corresponding frame back off. Every recursive algorithm can be rewritten as a loop that
manages its own explicit stack instead of relying on the language's call stack, which is
one reason the stack ADT and recursion keep showing up together.

\section{Traps people fall into}

\begin{itemize}[leftmargin=1.4em]
  \item \textbf{Popping or peeking an empty stack.} Always check isEmpty first, or catch
    the exception the operation throws. This is the single most common stack bug.
  \item \textbf{Off-by-one on the top index.} Decide once whether an array-based stack's
    top index points at the last used slot or at the next free slot, and be consistent.
    Mixing the two conventions across push and pop is a frequent source of bugs.
  \item \textbf{Forgetting to resize (or checking too late).} An array-based stack that
    does not check capacity before writing will overflow the array's bounds instead of
    reporting overflow cleanly.
  \item \textbf{Leaving a stale reference in Java.} After popping from an array-based
    stack, the old slot still refers to the popped object unless you explicitly clear it,
    which can keep that object alive longer than intended.
  \item \textbf{Forgetting to delete in C/C++.} Popping a node from a linked stack does not
    free its memory by itself. If you are not using smart pointers, you still need to call
    delete on the node you just unlinked, or the memory leaks.
\end{itemize}

\section{Where this lives in the textbooks}

Both books use the wall to name the ADT chapter, so it is easy to find.

\begin{itemize}[leftmargin=1.4em]
  \item \textbf{Java, 3rd edition (Prichard and Carrano).} Chapter 7, Stacks, covers the
    ADT itself (7.1), simple applications such as checking for balanced braces (7.2), three
    implementations, array-based, reference-based, and one built from the ADT list, plus a
    comparison and the Java Collections Framework's own Stack class (7.3), postfix
    evaluation and infix-to-postfix conversion (7.4), a search-problem application (7.5),
    and the relationship between stacks and recursion (7.6).
  \item \textbf{C/C++, 8th edition (Carrano and Henry).} Chapter 8, Stacks, covers the ADT,
    simple uses of a stack, using a stack to evaluate postfix expressions, and the
    relationship between stacks and recursion. Chapter 9, Stack Implementation, covers an
    array-based implementation, a link-based implementation, and a version built with
    exceptions.
\end{itemize}

\end{document}
